% !TeX program = xelatex
% Standalone document: the C listing is embedded, with no external inputs.
\documentclass[UTF8,fontset=fandol,zihao=-4,a4paper]{ctexart}
\usepackage[a4paper,margin=2.35cm,headheight=16pt]{geometry}
\usepackage{amsmath,amssymb,bm}
\usepackage{booktabs,array,longtable}
\usepackage{enumitem,fancyhdr,lastpage,xcolor,listings}
\usepackage{hyperref}
\newcommand{\coursename}{工程硕士数学}
\newcommand{\studentname}{陈天乐}
\newcommand{\studentid}{[student ID removed]}
\newcommand{\classinfo}{深微硕6班}
\newcommand{\majorinfo}{集成电路工程}
\newcommand{\reportdate}{2026年10月7日}
\newcommand{\psif}{\psi}
\newcommand{\dd}{\mathrm{d}}
\newcommand{\abs}[1]{\left|#1\right|}
\setlength{\parindent}{2em}
\setlength{\parskip}{0.3em}
\linespread{1.2}
\setlist[enumerate]{itemsep=0.2em,topsep=0.3em}
\setlist[itemize]{itemsep=0.2em,topsep=0.3em}
\pagestyle{fancy}
\fancyhf{}
\lhead{\coursename}
\rhead{实验一：Hamming 级数求和}
\cfoot{第 \thepage\ 页，共 \pageref{LastPage} 页}
\renewcommand{\headrulewidth}{0.4pt}
\hypersetup{colorlinks=true,linkcolor=blue!50!black,urlcolor=blue!60!black,
 pdftitle={工程硕士数学第一次实验报告：Hamming级数求和},pdfauthor={陈天乐}}
\lstdefinestyle{cprogram}{language=C,basicstyle=\ttfamily\footnotesize,
 keywordstyle=\color{blue!60!black},commentstyle=\color{green!40!black},
 stringstyle=\color{red!55!black},numbers=left,numberstyle=\tiny\color{gray},
 numbersep=7pt,frame=single,rulecolor=\color{gray!50},breaklines=true,
 columns=fullflexible,keepspaces=true,showstringspaces=false,
 tabsize=4,xleftmargin=1.6em,framexleftmargin=1.3em,
 aboveskip=0.7em,belowskip=0.7em}
\lstset{style=cprogram}
\ctexset{section={format=\Large\bfseries},subsection={format=\large\bfseries}}

\begin{document}
\begin{center}
 {\LARGE\bfseries \coursename\quad 第一次实验报告}\par
 \vspace{0.55em}
 {\Large\bfseries Hamming 级数求和}\par
 \vspace{1em}
 \begin{tabular}{rl@{\qquad}rl}
  专业：&\majorinfo & 班级：&\classinfo\\
  姓名：&\studentname & 学号：&\studentid\\
  日期：&\reportdate & 实验编号：&EXP1
 \end{tabular}
\end{center}
\vspace{0.5em}\hrule\vspace{0.6em}

\section{实验题目名称}
\textbf{实验一：Hamming 级数求和。}计算
\begin{equation}\label{eq:original}
 \psif(x)=\sum_{k=1}^{\infty}\frac{1}{k(k+x)},
 \qquad x=0,0.1,0.2,\ldots,300,
\end{equation}
共3001个函数值，要求每个结果的\textbf{绝对误差小于 $10^{-10}$}，并研究如何减少计算步骤和运行时间。输出按“$x$、$\psif(x)$”两列排列，函数值保留12位小数。

%\textbf{题目示例校正：}指导材料第一行写为“$0.10\quad1.644934066848$”，但后一个数对应 $x=0$，并非 $x=0.1$。正确值为
%\[
% \psif(0)=\frac{\pi^2}{6}=1.644934066848\cdots,\qquad
% \psif(0.1)=1.534607244905\cdots.
%\]
%本实验以题目规定的网格 $x=0,0.1,\ldots,300$ 为准，不沿用该处笔误。

\section{实验目的和意义}
\begin{enumerate}
 \item 区分截断误差、浮点舍入误差和输出舍入误差，建立能够检验的误差控制准则。
 \item 认识正项级数的收敛速度，说明“当前项很小”并不等于“剩余总和很小”。
 \item 利用函数的递推关系共享计算结果，将3001次独立求和转化为少量初值计算和线性复杂度的递推。
 \item 掌握 Euler--Maclaurin 尾项修正以及稳定计算小量差值的方法，体会算法设计比单纯提高硬件速度更重要。
 \item 使用 C 语言完成程序设计、编译、数值检查和计时，通过独立高精度计算验证全部输出。
\end{enumerate}

本方法的理论意义是将无限求和、解析尾项和稳定递推结合起来；实用价值是在保证精度的同时显著降低计算量。类似思想也适用于具有平移递推关系的特殊函数和数值表生成。

\section{计算公式与算法}
\subsection{直接截断求和及其误差分析}
记前N项和为
\[
 S_N(x)=\sum_{k=1}^{N}\frac{1}{k(k+x)},\qquad
 R_N(x)=\psif(x)-S_N(x)
\]

对 $x\ge0$，$f_x(t)=1/[t(t+x)]$ 是正的单调递减函数，因此由积分比较法可得
\begin{equation}\label{eq:integralbound}
 \int_{N+1}^{\infty} f_x(t)\,\dd t
 \le R_N(x)\le
 \int_N^{\infty} f_x(t)\,\dd t
\end{equation}

当 $x>0$ 时，
\begin{equation}
 \frac{\log(1+\frac{x}{N+1})}{x}
 \le R_N(x)\le
 \frac{\log(1+\frac{x}{N})}{x}
\end{equation}

当 $x=0$ 时，取极限得到
\begin{equation}\label{eq:zerotail}
 \frac{1}{N+1}\le R_N(0)\le\frac{1}{N}
\end{equation}

由于 $R_N(x)\le R_N(0)$，统一取 $N=10^{10}$ 可以将\textbf{精确算术下}的截断误差控制在 $10^{-10}$ 附近以下，但3001点总共需要约
\[
 3001\times10^{10}=3.001\times10^{13}
\]

次求和项计算，且尚未控制浮点舍入误差。这就是直接计算速度很慢的主要原因。

%此外，若仅以 $1/[k(k+x)]<10^{-10}$ 作为停止条件，在 $x=0$ 时约 $k=10^5$ 就会停止，而剩余总和仍约为 $10^{-5}$，与要求相差五个数量级。

\subsection{递推链分解}
对 $x>0$ 作部分分式分解：
\[
 \frac{1}{k(k+x)}=\frac{1}{x}\left(\frac1k-\frac1{k+x}\right).
\]
令 $F(x)=x\psif(x)$，并规定 $F(0)=0$。对收敛的差级数进行望远镜求和，得到
\begin{align*}
 F(x+1)-F(x)
 &=\sum_{k=1}^{\infty}\left(\frac1{k+x}-\frac1{k+x+1}\right)\\
 &=\frac1{x+1}.
\end{align*}
于是
\begin{equation}\label{eq:recurrence}
 \boxed{\psif(x+1)=\frac{x}{x+1}\psif(x)+\frac{1}{(x+1)^2}.}
\end{equation}
$\psif(0)$ 单独由 $\pi^2/6$ 计算；其余点分成10条链，起点分别为
\[
 0.1,0.2,\ldots,0.9,1.0.
\]
例如 $0.1,1.1,\ldots,299.1$ 属于同一条链，$1,2,\ldots,300$ 属于另一条链。每条链有300点，因此仅需10个初值和 $10\times299=2990$ 次递推。

设 $p_i=\psif(i/10)$，则适于程序实现的形式为
\begin{equation}\label{eq:indexrec}
 \boxed{p_i=\frac{i-10}{i}p_{i-10}+\frac{100}{i^2},
 \qquad i=11,\ldots,3000.}
\end{equation}

使用整数下标，而不是反复执行 $x\leftarrow x+0.1$，可以避免网格值累积偏移。

\subsection{递推稳定性及整体误差控制}
若初值误差为 $e_0$，在精确递推下，由式\eqref{eq:recurrence}有
\[
 e_{j+1}=\frac{x_0+j}{x_0+j+1}e_j,
 \qquad e_j=\frac{x_0}{x_0+j}e_0.
\]

因此初值误差随 $x$ 增大而衰减，\textbf{正向递推不放大初值误差}；反向递推则会乘上大于1的系数，不宜使用。

采用 IEEE 754 双精度时，单位舍入误差 $u=2^{-53}\approx1.11\times10^{-16}$。在本题范围内，整数 $i$ 和 $i^2$ 可精确表示；式\eqref{eq:indexrec}的两个加数均为正，不存在灾难性相消。令 $\gamma_3=3u/(1-3u)$，每步正数递推的局部舍入误差可按 $\gamma_3$ 量级估计。每条链最多299步，且 $\psif(x)<1.65$，递推舍入的保守量级为
\[
 1.65\times299\times\gamma_3<1.65\times10^{-13}.
\]

结合式\eqref{eq:bound}，若对初值浮点计算另保留 $10^{-14}$ 的预算，则总误差量级仍小于 $2.5\times10^{-13}$；保留12位小数带来的绝对舍入误差不超过 $5\times10^{-13}$，远小于实验限值。


\subsection{用 Euler--Maclaurin 公式计算10个初值}
在 $0<x\le1$ 上，先直接计算前 $N$ 项，再对无限尾项作 Euler--Maclaurin 修正。保留 $B_2,B_4,B_6,B_8$ 四个修正项，有
\begin{align}
 \psif(x)
 ={}&S_N(x)+\int_N^{\infty}f_x(t)\,\dd t-\frac{f_x(N)}2\notag\\
 &-\sum_{r=1}^{4}\frac{B_{2r}}{(2r)!}f_x^{(2r-1)}(N)
 +E_4(x,N),\label{eq:EMgeneral}
\end{align}
其中
\[
 B_2=\frac16,\quad B_4=-\frac1{30},\quad
 B_6=\frac1{42},\quad B_8=-\frac1{30}.
\]
由
\[
 f_x^{(j)}(t)=\frac{(-1)^j j!}{x}
 \left(t^{-j-1}-(t+x)^{-j-1}\right)
\]
和
\[
 \int_N^{\infty}f_x(t)\,\dd t=\frac{\log(1+x/N)}{x},
\]
可得实际使用的近似式
\begin{equation}\label{eq:EMexplicit}
 \boxed{\begin{aligned}
 A_N(x)={}&\sum_{k=1}^{N}\frac1{k(k+x)}
 +\frac{\log(1+x/N)}{x}-\frac1{2N(N+x)}\\
 &+\sum_{r=1}^{4}\frac{B_{2r}}{2r\,x}
 \left(N^{-2r}-(N+x)^{-2r}\right).
 \end{aligned}}
\end{equation}
本实验统一选择 $\boxed{N=20}$。公式仅用于 $0<x\le1$ 的10个初值，$x=0$ 不代入含 $1/x$ 的表达式。

\subsection{初值截断误差的严格上界}
Euler--Maclaurin 余项的周期 Bernoulli 函数表示给出
\[
 |E_4(x,N)|\le\frac{|B_8|}{8!}
 \int_N^{\infty}|f_x^{(8)}(t)|\,\dd t.
\]

因为 $f_x^{(8)}(t)>0$，且 $f_x^{(7)}(\infty)=0$，所以
\begin{align*}
 |E_4(x,N)|
 &\le \frac{|B_8|}{8x}\left(N^{-8}-(N+x)^{-8}\right)\\
 &\le |B_8|N^{-9}=\frac{1}{30N^9}.
\end{align*}

第二步使用中值定理 $N^{-8}-(N+x)^{-8}\le8xN^{-9}$。代入 $N=20$，得到统一上界
\begin{equation}\label{eq:bound}
 \boxed{|E_4(x,20)|\le\frac1{30\cdot20^9}
 =6.5104166667\times10^{-14}.}
\end{equation}

该上界适用于全部10个初值，明显小于 $10^{-10}$，为舍入误差和输出舍入保留了充足余量。

\subsection{避免尾项计算中的相消}
在式\eqref{eq:EMexplicit}中，直接计算 $\log(1+x/N)$ 或两个接近的幂的差，会损失部分有效数字。令 $L=\log(1+x/N)$，则
\begin{equation}
 N^{-2r}-(N+x)^{-2r}
 =N^{-2r}\bigl[-\operatorname{expm1}(-2rL)\bigr].
\end{equation}

程序用 \texttt{log1p(x/N)} 计算 $L$，用 \texttt{expm1} 计算 $e^z-1$，而不先形成 $1+x/N$ 或 $e^z$ 再相减。

%\textbf{说明：}上述截断误差界是严格的解析界；浮点估计依赖通常的双精度运算模型及近机器精度的数学库。C 标准不为 \texttt{log1p}、\texttt{expm1} 指定统一的误差界，因此本实验还对所有实际结果进行独立80位高精度校验，不把误差预算当作未经验证的跨平台保证。

\section{结构程序设计}
\subsection{程序模块与流程}
\begin{center}
\begin{tabular}{p{3.4cm}p{10.6cm}}
\toprule
模块 & 功能\\
\midrule
\texttt{base\_value} & 前20项直接求和、解析积分尾项及四个 Bernoulli 修正项，生成初值。\\
\texttt{compute\_all} & 计算 $\psif(0)$ 和10个初值，分10条链完成2990次递推。\\
\texttt{direct\_sum} & 普通正向直接求和，仅用于对照实验。\\
\texttt{now\_seconds} & Windows 下使用高分辨率计时器；其他平台采用标准 C 计时。\\
\texttt{main} & 检查结果、保存3001点、显示代表点、进行批量计时和对照试验。\\
\bottomrule
\end{tabular}
\end{center}

程序主流程为：
\begin{enumerate}
 \item 设置 $N=20$、修正阶数4及网格下标上限3000，使用 \texttt{double}。
 \item 计算 $p_0=\pi^2/6$。
 \item 对 $r=1,\ldots,10$，令 $x_0=r/10$，用式\eqref{eq:EMexplicit}求出 $p_r$。
 \item 对每个 $r$，依次令 $i=r+10,r+20,\ldots\le3000$，按式\eqref{eq:indexrec}更新 $p_i$。
 \item 检查全部结果是否有限、为正且严格递减；按下标顺序写出文件。
 \item 对同一可执行程序进行计时，并由独立高精度程序逐点核验。
\end{enumerate}



\subsection{数据结构、输出与复杂度}
使用长度3001的 \texttt{double} 数组保存结果，内存约为 $3001\times8=24008$ bytes，即23.45 KiB。算法实际计算200个直接求和项、40个尾项修正和2990个递推步骤。设网格点数为 $M$，初值直接求和长度为 $N$，则计算复杂度为
\[
 O(10N+M),\qquad \text{空间复杂度为 }O(M).
\]

若只需顺序输出，可只保存10条链的最新值，将工作存储降到常数级；本实验为方便检查和计时，保留全部数组。

程序生成以下文件：
\begin{itemize}
 \item \texttt{hamming\_results.txt}：3001行、两列，函数值保留12位小数，与题目输出形式一致。
 \item \texttt{hamming\_results.csv}：除12位小数外，另保存17位有效数字，供独立校验和数据处理使用。
 \item \texttt{benchmark.txt}：计时方式、计算次数、批量耗时和直接求和对照数据。
\end{itemize}
完整可编译 C 程序见附录，同时作为 \texttt{hamming.c} 独立保存。


高精度校验使用 Python 3.14 与 mpmath 1.3.0，十进制工作精度设为80位。独立参考表达式为
\begin{equation}\label{eq:reference}
 \psif(x)=\frac{\operatorname{digamma}(1+x)+\gamma}{x}\quad(x>0),
 \qquad \psif(0)=\zeta(2),
\end{equation}

其中 $\gamma$ 是 Euler 常数，$\operatorname{digamma}$ 是对数 Gamma 函数的导数。


%\subsection{编译与实际验证环境}
%本次运行环境为 Windows x64，编译器为 Microsoft Visual C 14.44.35207，采用优化设置 \texttt{/O2 /std:c11 /fp:precise}，未启用不保留标准浮点语义的快速数学优化。程序已编译并运行，所有3001个结果均由此 C 程序生成，而非仅列出理论值。

%高精度校验使用 Python 3.14 与 mpmath 1.3.0，十进制工作精度设为80位。独立参考表达式为
%\begin{equation}\label{eq:reference}
% \psif(x)=\frac{\operatorname{digamma}(1+x)+\gamma}{x}\quad(x>0),
% \qquad \psif(0)=\zeta(2),
%\end{equation}
%其中 $\gamma$ 是 Euler 常数，$\operatorname{digamma}$ 是对数 Gamma 函数的导数。此式仅用于\textbf{独立验证}，未在 C 求值算法中调用。参考网格用精确有理数 $i/10$ 构造，以免把二进制网格近似当成真值。

\section{结果讨论和分析}
\subsection{代表性计算结果}
3001个结果全部已保存，下面列出代表点。
\begin{center}
\begin{tabular}{r@{\qquad}r|r@{\qquad}r}
\toprule
$x$ & $\psif(x)$ & $x$ & $\psif(x)$\\
\midrule
0.00 & 1.644934066848 & 10.00 & 0.292896825397\\
0.10 & 1.534607244905 & 20.00 & 0.179886982857\\
0.20 & 1.440878841547 & 50.00 & 0.089984106767\\
0.50 & 1.227411277760 & 100.00 & 0.051873775176\\
1.00 & 1.000000000000 & 150.00 & 0.037274537258\\
2.00 & 0.750000000000 & 200.00 & 0.029390154741\\
5.00 & 0.456666666667 & 250.00 & 0.024402700998\\
\multicolumn{2}{c|}{} & 300.00 & 0.020942212934\\
\bottomrule
\end{tabular}
\end{center}

当 $x=m\in\mathbb{N}_{+}$ 时，可由望远镜求和得到
\[
 \psif(m)=\frac{H_m}{m},\qquad H_m=\sum_{k=1}^{m}\frac1k.
\]

因此 $\psif(1)=1$、$\psif(2)=3/4$，与程序输出一致。全部值为正，且严格递减，这也与
\[
 \psif'(x)=-\sum_{k=1}^{\infty}\frac1{k(k+x)^2}<0
\]

相符。特别地，$x=300$ 的结果是 $0.020942212934$，而并非直接用粗略的 $\log(x)/x$ 代替。

\subsection{全部3001点的精度核验}
以式\eqref{eq:reference}的80位计算结果为参考，分别比较 C 内部双精度值和实际写出的12位小数值，得到：
\begin{center}
\begin{tabular}{p{7.3cm}r@{\quad}l}
\toprule
校验项目 & 数值 & 位置或说明\\
\midrule
初值解析截断误差统一上界 & $6.51042\times10^{-14}$ & $0<x\le1$\\
10个初值的最大实际解析截断误差 & $3.56980\times10^{-16}$ & $x=0.1$\\
全部内部值的最大绝对误差 & $5.23537\times10^{-16}$ & $x=0.9$\\
12位小数输出的最大绝对误差 & $4.99948\times10^{-13}$ & $x=160.8$\\
满足绝对误差小于 $10^{-10}$ 的点数 & 3001 & 全部通过\\
\bottomrule
\end{tabular}
\end{center}

这里的“实际解析截断误差”是以80位算术计算 $A_{20}(x)$ 再与参考值比较，因而不与 C 程序舍入误差混淆。内部值的最大误差约为限值的 $1/191000$；最终文本输出的最大误差约为限值的 $1/200$。

\subsection{直接求和对照试验}
取 $x=0$，直接用 \texttt{double} 从 $k=1$ 向上累加；误差仍使用80位参考值计算。实际结果如下。
\begin{center}
\small
\begin{tabular}{r r r r}
\toprule
截断项数 $N$ & 直接和（约15位小数） & 实际绝对误差 & 单点耗时（秒）\\
\midrule
$10^3$ & 1.643934566681562 & $9.99500\times10^{-4}$ & $7.64080\times10^{-7}$\\
$10^4$ & 1.644834071848065 & $9.99950\times10^{-5}$ & $7.63094\times10^{-6}$\\
$10^5$ & 1.644924066898242 & $9.99995\times10^{-6}$ & $7.65704\times10^{-5}$\\
$10^6$ & 1.644933066848770 & $9.99999\times10^{-7}$ & $7.72392\times10^{-4}$\\
$10^8$ & 1.644934057834575 & $9.01365\times10^{-9}$ & $7.59942\times10^{-2}$\\
$2\times10^8$ & 1.644934057834575 & $9.01365\times10^{-9}$ & $1.53569\times10^{-1}$\\
\bottomrule
\end{tabular}
\end{center}

前四行符合 $R_N(0)\approx1/N$。最后两行直接和完全相同，尽管计算量增加一倍。

其原因是，当部分和约为1.64时，相邻双精度浮点数的间隔约为 $2.22\times10^{-16}$。一项 $1/k^2$ 小于半个间隔后，加到部分和上会被舍掉，发生于
\[
 k\gtrsim\left(1.11\times10^{-16}\right)^{-1/2}
 \approx9.49\times10^7.
\]

因此，普通从大项到小项的累加不仅速度慢，后期还会停滞，不能靠无限增加 $N$ 来保证 $10^{-10}$ 精度。倒序求和、成对求和或 Kahan 补偿求和可减轻舍入损失，但\textbf{不能消除原级数 $1/N$ 级别的截断尾项}，仍需尾项修正或其他加速算法。

\subsection{运行时间与计算量比较}
加速程序重复计算整张3001点表10000次，测得总耗时为 $0.074388900$ 秒，平均每张表
\[
 \boxed{7.43889\times10^{-6}\ \text{秒，约 }7.44\ \text{微秒}.}
\]

计时采用 Windows 的 \texttt{QueryPerformanceCounter}，不包含文件写入和高精度校验；为减少短时测量误差使用重复批量计时，并以 \texttt{volatile} 读取和累加结果防止无用计算被优化掉。对照试验前三类及 $10^6$ 项试验也采用重复计时，$10^8$ 与 $2\times10^8$ 项各测一次。计时反映本机本次运行，不作为跨机器的固定性能指标。

\begin{center}
\begin{tabular}{p{4.1cm}p{4.6cm}p{5.2cm}}
\toprule
方法 & 计算量/每张表 & 精度与特点\\
\midrule
直接截断到 $10^6$ 项 & 约 $3.001\times10^9$ 项 & 截断误差约 $10^{-6}$，不达标。\\
直接截断到 $10^{10}$ 项 & 约 $3.001\times10^{13}$ 项 & 仅在精确算术下控制尾项；普通双精度累加会停滞。\\
尾项修正与正向递推 & 200项、40个修正项、2990步 & 3001点实际输出误差均小于 $10^{-10}$。\\
\bottomrule
\end{tabular}
\end{center}

按本机 $N=10^6$ 的\emph{单点}计时外推，3001点独立直接求和需约 $2.32$ 秒，而加速方法计算整表约 $7.44$ 微秒。该约 $3.1\times10^5$ 的比值只是操作量相近时的\textbf{估计}，并非实测整表直接求和的加速比；且对照的 $10^6$ 项方法尚未达标，不能将两者视为同精度算法。真正重要的是加速方法既降低了计算量，又具有明确的尾项误差界。

\subsection{参数、算法特点及可改进之处}
\begin{enumerate}
 \item \textbf{初值精度：}初值误差沿递推链乘上 $x_0/x$，不会放大，但较小的 $x$ 最接近初值精度，因此应优先控制10个初值的误差。
 \item \textbf{截断长度：}四个修正项的统一余项界为 $1/(30N^9)$。例如 $N=10$ 时为 $3.33\times10^{-11}$，理论上已接近达标；选择 $N=20$ 可显著增加舍入余量，而只多计算100个求和项。
 \item \textbf{递推方向：}采用正向递推。反向递推会放大误差，且靠近 $x=0$ 时还需额外处理。
 \item \textbf{准确性与成本：}不必对每个 $x$ 独立调用高精度函数；仅在实验验证阶段用高精度参考值检查程序，不把验证成本计入求值速度。
 \item \textbf{进一步优化：}可直接用 $\psif(1)=1$ 替代第10个初值求和；若只需流式输出，可用10个工作值而不是完整数组；也可按目标精度自动选择 $N$ 与修正阶数。
 \item \textbf{应用边界：}本算法利用步长0.1网格中相差1的点共享递推。若改为无规律的任意采样点，需要重新组织初值与递推链，而不是机械复用当前索引关系。
\end{enumerate}

\subsection{程序实现和调试中的问题及处理}
\begin{enumerate}
 \item \textbf{$x=0$ 的除零问题：}初值尾项公式含 $1/x$，不能直接代入0。将 $\psif(0)=\pi^2/6$ 单独计算。
 %\item \textbf{示例首行不一致：}通过解析特殊值和高精度计算确认，$1.644934066848$ 应对应 $x=0$，输出采用正确网格。
 \item \textbf{网格漂移：}用整数索引 $i$ 生成 $i/10$，递推也直接使用索引公式，避免连续加0.1造成累计偏移。
 \item \textbf{相消及输出误差：}尾项中使用 \texttt{log1p}/\texttt{expm1}；既校验17位有效数字保存的内部值，也校验最终12位小数文本，不以显示小数位数代替误差证明。
 \item \textbf{计时被优化：}初次重复计时中，直接求和的相同输入调用被优化为重复使用已知结果，使时间读数接近零。改用 \texttt{volatile} 输入并保留计时累加结果后，得到随 $N$ 增长而近似线性增加的实际耗时。
 \item \textbf{大 $N$ 累加停滞：}通过 $10^8$、$2\times10^8$ 项对照确认舍入造成停滞。解决办法不是继续增加项数，而是更换为解析尾项修正与稳定递推。
\end{enumerate}

\subsection{实验结论}
本实验采用“\textbf{20项直接求和＋四个 Euler--Maclaurin 尾项修正＋10条正向递推链}”计算 Hamming 级数。初值解析截断误差统一不超过 $6.51042\times10^{-14}$，全网格实际内部最大绝对误差为 $5.23537\times10^{-16}$，12位小数输出的最大绝对误差为 $4.99948\times10^{-13}$，3001点全部满足 $10^{-10}$ 的精度要求。

直接求和之所以效率低难以计算，原因在于尾项仅按 $1/N$ 的速率衰减，并且3001个点有大量重复计算，以及普通浮点累加后期停滞。实验表明，利用解析关系得到递推式、明确控制尾项并选取稳定计算方式，可以同时改善精度与效率。

\clearpage
\appendix
\section{源代码}
\begin{lstlisting}
/* Hamming series: 3001 values on x = i/10, absolute target 1e-10.
   C99, double precision, no external libraries beyond the C math library. */
#if defined(_MSC_VER)
#define _CRT_SECURE_NO_WARNINGS
#endif
#include <stdio.h>
#include <stdlib.h>
#include <math.h>
#include <time.h>
#if defined(_WIN32)
#include <windows.h>
#endif

#define LAST_INDEX 3000
#define N_BASE 20
#define EM_ORDER 4
#define FAST_REPEATS 10000
#if defined(_MSC_VER)
#define NOINLINE __declspec(noinline)
#elif defined(__GNUC__) || defined(__clang__)
#define NOINLINE __attribute__((noinline))
#else
#define NOINLINE
#endif

static volatile double benchmark_sink = 0.0;
static volatile double direct_benchmark_x = 0.0;

static double now_seconds(void)
{
#if defined(_WIN32)
    LARGE_INTEGER ticks, frequency;
    QueryPerformanceCounter(&ticks);
    QueryPerformanceFrequency(&frequency);
    return (double)ticks.QuadPart/(double)frequency.QuadPart;
#else
    return (double)clock()/(double)CLOCKS_PER_SEC;
#endif
}

/* First N terms, followed by an Euler-Maclaurin tail. */
static double base_value(double x)
{
    const double B[EM_ORDER] = {1.0/6.0, -1.0/30.0,
                               1.0/42.0, -1.0/30.0};
    const double n = (double)N_BASE;
    const double l = log1p(x / n);
    double sum = 0.0, correction = 0.0;
    int k, r;
    for (k = 1; k <= N_BASE; ++k)
        sum += 1.0 / ((double)k * ((double)k + x));
    for (r = 1; r <= EM_ORDER; ++r) {
        double diff = pow(n, -2.0*r) * (-expm1(-2.0*r*l));
        correction += B[r-1] * diff / (2.0*r*x);
    }
    return sum + log1p(x/n)/x - 1.0/(2.0*n*(n+x))
               + correction;
}

/* Ten residue classes: exact integer indices avoid drifting x += 0.1. */
static NOINLINE void compute_all(double values[LAST_INDEX + 1])
{
    const double pi = acos(-1.0);
    int r, i;
    values[0] = pi*pi/6.0;
    for (r = 1; r <= 10; ++r) {
        values[r] = base_value((double)r/10.0);
        for (i = r+10; i <= LAST_INDEX; i += 10)
            values[i] = ((double)(i-10)/(double)i)*values[i-10]
                      + 100.0/((double)i*(double)i);
    }
}

static NOINLINE double direct_sum(double x, int n)
{
    double sum = 0.0;
    int k;
    for (k = 1; k <= n; ++k)
        sum += 1.0/((double)k*((double)k+x));
    return sum;
}

int main(void)
{
    double values[LAST_INDEX + 1];
    const int selected[] = {0,1,2,5,10,20,50,100,200,500,
                            1000,1500,2000,2500,3000};
    const int ns[] = {1000,10000,100000,1000000,100000000,200000000};
    FILE *txt, *csv, *bench;
    double begin, end;
    double seconds;
    int i, j, repeat;
    compute_all(values);
    for (i = 0; i <= LAST_INDEX; ++i) {
        if (!isfinite(values[i]) || values[i] <= 0.0 ||
            (i > 0 && values[i] >= values[i-1])) {
            fprintf(stderr, "Invalid or nondecreasing result at index %d\n", i);
            return EXIT_FAILURE;
        }
    }
    txt = fopen("hamming_results.txt", "w");
    csv = fopen("hamming_results.csv", "w");
    bench = fopen("benchmark.txt", "w");
    if (!txt || !csv || !bench) {
        fprintf(stderr, "Cannot create output files.\n");
        if (txt) fclose(txt);
        if (csv) fclose(csv);
        if (bench) fclose(bench);
        return EXIT_FAILURE;
    }
    fprintf(csv, "x,psi_12dp,psi_full_precision\n");
    for (i = 0; i <= LAST_INDEX; ++i) {
        fprintf(txt, "%6.2f  %.12f\n", (double)i/10.0, values[i]);
        fprintf(csv, "%.1f,%.12f,%.17g\n", (double)i/10.0,
                values[i], values[i]);
    }
    printf("Selected results (all 3001 values saved in TXT and CSV):\n");
    for (j = 0; j < (int)(sizeof(selected)/sizeof(selected[0])); ++j) {
        i = selected[j];
        printf("%6.2f  %.12f\n", (double)i/10.0, values[i]);
    }
    fprintf(bench, "algorithm=Euler-Maclaurin + forward recurrence\n");
#if defined(_WIN32)
    fprintf(bench, "timer=Windows QueryPerformanceCounter (elapsed time)\n");
#else
    fprintf(bench, "timer=C clock (CPU time)\n");
#endif
    fprintf(bench, "N=%d\nEM_order=%d\nbase_values=10\n", N_BASE, EM_ORDER);
    fprintf(bench, "base_term_evaluations=%d\nrecurrence_steps=2990\n",
            10*N_BASE);
    fprintf(bench, "truncation_bound=%.17g\n", 1.0/(30.0*pow((double)N_BASE,9.0)));
    begin = now_seconds();
    for (repeat = 0; repeat < FAST_REPEATS; ++repeat) {
        compute_all(values);
        benchmark_sink += values[repeat % (LAST_INDEX+1)];
    }
    end = now_seconds();
    seconds = end-begin;
    fprintf(bench, "fast_repeats=%d\nfast_batch_seconds=%.9f\n",
            FAST_REPEATS, seconds);
    fprintf(bench, "fast_seconds_per_3001_values=%.12g\n",
            seconds/(double)FAST_REPEATS);
    fprintf(bench, "direct_x=0\nN,sum_17digits,repeats,total_seconds,seconds_per_value\n");
    for (j = 0; j < (int)(sizeof(ns)/sizeof(ns[0])); ++j) {
        int n = ns[j];
        int repetitions = 100000000/n;
        if (repetitions < 1) repetitions = 1;
        double s = direct_sum(0.0, n);
        begin = now_seconds();
        for (repeat = 0; repeat < repetitions; ++repeat)
            /* Volatile input prevents the optimizer hoisting a pure call. */
            benchmark_sink += direct_sum(direct_benchmark_x, n);
        end = now_seconds();
        seconds = end-begin;
        fprintf(bench, "%d,%.17g,%d,%.9f,%.12g\n", n,s,repetitions,
                seconds,seconds/(double)repetitions);
    }
    i = (fclose(txt) != 0);
    i |= (fclose(csv) != 0);
    i |= (fclose(bench) != 0);
    if (i) {
        fprintf(stderr, "Output write failure.\n");
        return EXIT_FAILURE;
    }
    printf("Checks passed: finite, positive and strictly decreasing.\n");
    printf("Benchmark saved in benchmark.txt (excludes file output).\n");
    return EXIT_SUCCESS;
}
\end{lstlisting}

\section{附带文件清单}
\begin{itemize}
 \item \texttt{EXP1.tex}：本报告的独立 LaTeX 源文件，程序已内嵌，编译不依赖外部代码文件。
 \item \texttt{hamming.c}、\texttt{hamming.exe}：完整 C 源码及本机生成的 Windows 可执行程序。
 \item \texttt{hamming\_results.txt}、\texttt{hamming\_results.csv}：全部3001点的正式结果。
 \item \texttt{benchmark.txt}：本次运行的原始计时记录。
 \item \texttt{verify\_results.py}：独立80位高精度校验程序；需要 mpmath。
 \item \texttt{validation\_summary.json} 与逐点误差表：保存在“验证记录”文件夹，记录全部点的参考值和误差。
\end{itemize}
\end{document}
